% SPDX-License-Identifier: CC-BY-SA-4.0
% Author: Matthieu Perrin
% Part: <Nom de la partie>
% Section: <Nom de la section>
% Sub-section: <Nom de la sous-section>  % (facultatif, laisser vide si non utilisé)
% Frame: <Titre de la slide>

\begingroup

\begin{frame}{Problème de Correspondance de Post Initialisé Borné}

  \on[text, top=-2mm]{
    Soit $\Sigma$ un alphabet contenant au moins deux symboles. 
  }
  
  \on[text, top=0mm, width=1.15\textwidth]{
    \Probleme{bpcpi}{
      \begin{itemize}
      \item Un ensemble fini de \structure{sortes de dominos} {\small$D \in \mathscr{P}(\Sigma^+ \times \Sigma^+)$}
        \begin{itemize}
        \item Dominos de la forme \VDomino{$\alpha$}{$\beta$}, avec $\alpha$ et $\beta$ des mots sur $\Sigma$
        \end{itemize}
      \item Un \structure{domino initial} $d_0 = \VDomino{\gamma}{\delta} \in D$
      \item Une \structure{borne $b \in \mathbb{N}$} encodée en \alert{unaire} 
      \end{itemize}
    }{
      Existe-t-il un mot\footnote[frame]{On peut utiliser chaque domino jusqu'à $b$ fois}
      $\VDomino{\alpha_1}{\beta_1} \,\cdots\, \VDomino{\alpha_k}{\beta_k} \in D^+$ tel que :

      \vspace{1mm}
      \begin{minipage}{.27\textwidth}
        \begin{itemize}
        \item $\VDomino{\alpha_1}{\beta_1} = \VDomino{\gamma}{\delta}$
        \end{itemize}
      \end{minipage}%
      \begin{minipage}{.3\textwidth}
        \begin{itemize}
        \item $\structure{\alpha_{1} \cdots \alpha_{k}} = \example{\beta_{1} \cdots \beta_{k}}$
        \end{itemize}
      \end{minipage}
      \begin{minipage}{.2\textwidth}
        \begin{itemize}
        \item \alert{$k \le b$}
        \end{itemize}
      \end{minipage}
    }
  }

    \onExampleBlock[bottom=3mm] {Exemple} {
    \begin{itemize}
    \item $\left\langle\, \left\{
      \VDomino{a}{ab}, 
      \VDomino{baba}{a}, 
      \VDomino{bb}{ba}
      \right\},\,
      \VDomino{a}{ab},\, 1111 \,\right\rangle
      \quad\alert{\in\textsc{neg}(\textsc{bpcpi})}
      $
    \end{itemize}
  }

\end{frame}

\endgroup
